% Chapter: Formal Semantics of C⁴
% New chapter: categorical presheaf/sheaf model, denotation of all terms,
% soundness theorem, and correspondence with the operational semantics.

\chapter{Formal Semantics of \(\mathrm{C}^4\)}
\label{ch:c4-formal-semantics}

The typing rules of Chapter~\ref{ch:c4-calculus} give the
\emph{syntactic} face of $\mathrm{C}^4$.  This chapter constructs its
\emph{semantic} face: a precise mathematical model in which every
well-typed term denotes an actual mathematical object, and every
computation rule corresponds to an equality in the model.  The
semantic model simultaneously justifies the typing rules (soundness)
and the computation rules (adequacy), and it provides the mathematical
foundation for the Lean~4 mechanisation in Chapter~\ref{ch:c4-lean}.


% ─────────────────────────────────────────────────────────
\section{The Presheaf Model}
\label{sec:presheaf-model}

\subsection{The Base Category}

Fix a base type \(\alpha : \mathsf{Set}\) (the type of Python program inputs).
The \textbf{refinement site category} \(\mathcal{C}_\alpha\) is:
\begin{itemize}
\item \textbf{Objects:} Z3-decidable predicates
  \(\varphi : \alpha \to \mathsf{Prop}\), ordered by implication.
\item \textbf{Morphisms:}
  \(\mathrm{Hom}(\varphi, \psi) = \{\star \mid \forall x.\, \psi(x) \Rightarrow \varphi(x)\}\),
  i.e., morphisms go from stronger predicates to weaker ones
  (contravariant w.r.t.\ implication).
\item \textbf{Terminal object:} the predicate \(\mathsf{True}\).
\item \textbf{Products:}
  \(\varphi \times \psi = \varphi \wedge \psi\) (conjunction).
\end{itemize}

Note the contravariance: a morphism from \(\varphi\) to \(\psi\) in
\(\mathcal{C}_\alpha\) is a proof that \(\psi \Rightarrow \varphi\)
(restriction goes from big open sets to smaller ones).

\subsection{The Category of Presheaves}

\begin{definition}[Presheaf over the Refinement Site]
A \textbf{presheaf} on \(\mathcal{C}_\alpha\) is a functor
\(\mathcal{F} : \mathcal{C}_\alpha \to \mathsf{Set}\), i.e.\ a
family of sets \(\mathcal{F}(\varphi)\) for each predicate
\(\varphi\), together with restriction maps
\[
  \mathcal{F}(f) : \mathcal{F}(\psi) \to \mathcal{F}(\varphi)
\]
for each morphism \(f : \varphi \to \psi\) (i.e.\
\(\psi \Rightarrow \varphi\)), satisfying:
\[
  \mathcal{F}(\mathrm{id}) = \mathrm{id},
  \qquad
  \mathcal{F}(g \circ f) = \mathcal{F}(f) \circ \mathcal{F}(g).
\]
\end{definition}

We write \(\restr{s}{f}\) for the image of \(s \in \mathcal{F}(\psi)\)
under \(\mathcal{F}(f)\).

\subsection{The Sheaf Condition}

\begin{definition}[Sheaf over the Refinement Site]
A presheaf \(\mathcal{F}\) is a \textbf{sheaf} if for every cover
\(\covers = \{\varphi_i\}_{i \in I}\) of \(\mathsf{True}\)
(i.e.\ \(\forall x.\,\bigvee_i \varphi_i(x)\)), the following
diagram is an equalizer:
\[
  \mathcal{F}(\mathsf{True}) \;\xrightarrow{\;\delta\;}
  \prod_i \mathcal{F}(\varphi_i) \;\rightrightarrows\;
  \prod_{i,j} \mathcal{F}(\varphi_i \wedge \varphi_j)
\]
where \(\delta(s) = (\restr{s}{\varphi_i})_i\) and the two arrows
are \((s_i)_i \mapsto (\restr{s_i}{\varphi_i \wedge \varphi_j})_{i,j}\)
and \((s_i)_i \mapsto (\restr{s_j}{\varphi_i \wedge \varphi_j})_{i,j}\).
\end{definition}

In concrete terms: \(\mathcal{F}\) is a sheaf iff for every compatible
family of local sections \((s_i)_i\) (satisfying
\(\restr{s_i}{\varphi_i \wedge \varphi_j} = \restr{s_j}{\varphi_i
\wedge \varphi_j}\)), there exists a unique global section
\(s \in \mathcal{F}(\mathsf{True})\) restricting to each \(s_i\).

\begin{lemma}[Propositional Sheaves are Trivial]
\label{lem:prop-sheaf-trivial}
If \(\mathcal{F}\) is a presheaf with propositional values
(i.e.\ each \(\mathcal{F}(\varphi)\) is a proposition), then the
sheaf condition holds automatically.
\end{lemma}

\begin{proof}
The equalizer condition asks for a unique global section.
For propositions, all sections of the same proposition are equal
(proof irrelevance: \(\forall p, q : P.\, p = q\) in \(\mathsf{Prop}\)).
So uniqueness is automatic.  Existence follows from exhaustiveness:
choose any \(\varphi_k\) with \(\varphi_k(x)\) for the given
input \(x\), and take \(\restr{s_k}{x}\).
\end{proof}

This lemma is the semantic counterpart of the Hedberg collapse:
for verification purposes (where all proof goals are propositions),
the sheaf condition is vacuously satisfied.


% ─────────────────────────────────────────────────────────
\section{Interpretation of Types}
\label{sec:type-interpretation}

We define the \textbf{denotation} \(\den{\Gamma \vdash A}\) of each
well-formed type \(A\) in context \(\Gamma\) as a sheaf over
\(\mathcal{C}_\alpha\).

\subsection{Contexts}

A context \(\Gamma = (x_1 : A_1, \ldots, x_n : A_n)\) denotes a
sheaf of tuples:
\[
  \den{\Gamma}(\varphi) \;\coloneqq\;
  \{(a_1, \ldots, a_n) \mid
    a_1 \in \den{A_1}(\varphi),\;
    \ldots,\;
    a_n \in \den{A_n}(\varphi) \}.
\]

A fiber context \(\Gamma, x :_\psi A\) adds the constraint
\(\psi(x)\):
\[
  \den{\Gamma, x :_\psi A}(\varphi) \;\coloneqq\;
  \den{\Gamma}(\varphi) \times
  \{a \in \den{A}(\varphi \wedge \psi) \mid \text{well-typed}\}.
\]

\subsection{Universe Levels}

Universes are interpreted as sets of sheaves:
\[
  \den{\mathcal{U}_i^{\site}}(\varphi)
  \;\coloneqq\;
  \{\text{sheaves over } \restr{\site}{\varphi}\text{ at level } i\}.
\]

\subsection{Type Denotations}

\begin{align*}
  \den{\Pi(x : A).\, B}(\varphi)
  &\;\coloneqq\;
  \{f \mid \forall \psi \leq \varphi,\; \forall a \in \den{A}(\psi).\;
    f(\psi, a) \in \den{B}(\psi)[x \mapsto a]\} \\[4pt]
  \den{\Sigma(x : A).\, B}(\varphi)
  &\;\coloneqq\;
  \{(a, b) \mid a \in \den{A}(\varphi),\;
    b \in \den{B}(\varphi)[x \mapsto a]\} \\[4pt]
  \den{\Path_A\, a\, b}(\varphi)
  &\;\coloneqq\;
  \{\gamma : \interval \to \den{A}(\varphi) \mid
    \gamma(0) = \den{a}(\varphi),\; \gamma(1) = \den{b}(\varphi)\} \\[4pt]
  \den{\{x : A \mid \psi(x)\}}(\varphi)
  &\;\coloneqq\;
  \{a \in \den{A}(\varphi \wedge \psi)\} \\[4pt]
  \den{\mathsf{Ind}(\Delta)}(\varphi)
  &\;\coloneqq\;
  \text{initial algebra of the signature } \Delta \text{ in }
  \mathsf{Set} \\[4pt]
  \den{\mathcal{U}_\PyObj}(\varphi)
  &\;\coloneqq\;
  \mathsf{OTerm}_\varphi
  \quad\text{(OTerms modulo axioms restricted to } \varphi\text{)}
\end{align*}

\begin{proposition}[Well-definedness of type denotation]
For each well-formed type \(\Gamma \vdash A : \mathcal{U}_i\),
the denotation \(\den{A} : \mathcal{C}_\alpha \to \mathsf{Set}\)
is a sheaf (not just a presheaf).
\end{proposition}

\begin{proof}
By structural induction on \(A\):
\begin{itemize}
\item \(\Pi\): The sheaf condition for function sheaves follows from
  the sheaf condition on \(B\) (applying it pointwise on \(A\)-valued
  arguments).
\item \(\Sigma\): Follows from the product of sheaves being a sheaf.
\item \(\Path\): The path sheaf condition follows from the fact that
  path spaces in sets are sets.
\item Refinement: Lemma~\ref{lem:prop-sheaf-trivial} applies.
\item \(\mathrm{Ind}\): The initial algebra is a sheaf because
  it has decidable equality (for the base types used in Python).
\item \(\mathcal{U}_\PyObj\): OTerms form a sheaf because the
  CCTT axiom set is finite and decidable.
\qed
\end{itemize}
\end{proof}


% ─────────────────────────────────────────────────────────
\section{Interpretation of Terms}
\label{sec:term-interpretation}

For each well-typed term \(\Gamma \vdash t : A\), we define a
\textbf{global section} \(\den{t} : \prod_\varphi \den{A}(\varphi)\)
that is natural in \(\varphi\) (i.e., a section of the sheaf
\(\den{A}\) over the terminal object).

\subsection{Variables and Abstraction}

\[
  \den{x}(\varphi, \rho)
  \;\coloneqq\; \rho(x)
  \quad\text{(lookup in the environment \(\rho : \den{\Gamma}(\varphi)\))}
\]

\[
  \den{\lambda(x : A).\, t}(\varphi, \rho)
  \;\coloneqq\;
  \psi \mapsto a \mapsto \den{t}(\psi, \rho[\psi] \cup \{x \mapsto a\})
\]
where \(\rho[\psi]\) restricts the environment to the sub-fiber \(\psi\).

\subsection{Path Terms}

\[
  \den{\langle i \rangle\, t}(\varphi, \rho)
  \;\coloneqq\;
  r \mapsto \den{t}(\varphi, \rho[i \mapsto r])
\]

\[
  \den{p\; r}(\varphi, \rho)
  \;\coloneqq\; \den{p}(\varphi, \rho)(\den{r})
\]

\subsection{Transport}

\[
  \den{\transp^i\, A\, h\, t}(\varphi, \rho)
  \;\coloneqq\;
  \mathrm{transp}_{\den{A}}\bigl(\den{t}(\varphi, \rho)\bigr)
\]
where \(\mathrm{transp}_F\) is the transport map of the fiber functor
\(F = \den{A} : [0, 1] \to \mathsf{Set}\) evaluated at the path
\(\den{h}\).

\subsection{Homogeneous Composition}

\[
  \den{\hfilling^i\, A\, [\psi_k \mapsto u_k]\, t_0}(\varphi, \rho)
  \;\coloneqq\;
  \mathrm{fill}\bigl(
  \{\den{\psi_k}\}, \{\den{u_k}(\varphi, \rho)\}, \den{t_0}(\varphi, \rho)
  \bigr)
\]
where \(\mathrm{fill}\) is the unique global section guaranteed by the
sheaf condition applied to the partial boundary data
\(\{\den{u_k}\}\).

\subsection{Descent (The Key Denotation)}

\begin{definition}[Denotation of Descent]
\label{def:descent-denotation}
For a descent term
\(\mathsf{fill}(\covers, \{t_\alpha\}, \{p_{\alpha\beta}\})\)
where \(\covers = \{\varphi_i\}_{i \in I}\):
\[
  \den{\mathsf{fill}(\covers, \{t_\alpha\}, \{p_{\alpha\beta}\})}
  (\mathsf{True}, \rho)
  \;\coloneqq\;
  \text{the unique } s \in \den{A}(\mathsf{True})
  \text{ with }
  \restr{s}{\varphi_\alpha} = \den{t_\alpha}(\varphi_\alpha, \rho)
\]
The uniqueness follows from the sheaf condition on \(\den{A}\).
The compatibility hypothesis
\(p_{\alpha\beta} : \restr{t_\alpha}{\varphi_\alpha \wedge \varphi_\beta}
\equiv \restr{t_\beta}{\varphi_\alpha \wedge \varphi_\beta}\)
ensures that the local data is consistent.
\end{definition}

\begin{proposition}[Descent denotation is well-defined]
The denotation \(\den{\mathsf{fill}(\covers, \{t_\alpha\}, \{p_{\alpha\beta}\})}\)
exists and is a global section of \(\den{A}\).
\end{proposition}

\begin{proof}
The type \(\den{A}\) is a sheaf (by the preceding proposition).
The local sections \(\den{t_\alpha}\) are compatible on overlaps
(by the denotation of the compatibility proofs \(p_{\alpha\beta}\)).
The sheaf condition gives a unique global section.
For propositional \(\den{A}\), Lemma~\ref{lem:prop-sheaf-trivial}
applies directly.
\end{proof}

\subsection{Restriction}

\[
  \den{\restr{t}{\varphi}}(\psi, \rho)
  \;\coloneqq\;
  \den{t}(\psi \wedge \varphi, \rho)
\]

\subsection{OTerm Embedding}

\[
  \den{\den{P}}(\varphi, \rho)
  \;\coloneqq\;
  \llbracket P \rrbracket_{\mathsf{CCTT}, \varphi}
\]
the CCTT denotation of the Python AST \(P\) restricted to inputs
satisfying \(\varphi\).


% ─────────────────────────────────────────────────────────
\section{The Soundness Theorem}
\label{sec:soundness}

\begin{theorem}[Semantic Soundness of \(\mathrm{C}^4\)]
\label{thm:soundness}
Let \(\Gamma \vdash t : A\) be a valid typing judgment.  Then:
\begin{enumerate}
\item \(\den{A}\) is a sheaf over \(\mathcal{C}_\alpha\).
\item \(\den{t}\) is a global section of \(\den{A}\):
  \[
    \den{t} \;\in\; \mathcal{F}(\mathsf{True})
    \quad\text{where } \mathcal{F} = \den{A}.
  \]
\item The denotation is \textbf{natural}: for every site morphism
  \(f : \psi \to \varphi\) (i.e.\ \(\varphi \Rightarrow \psi\)),
  \[
    \restr{\den{t}}{\psi} \;=\; \den{\restr{t}{\psi}}.
  \]
\end{enumerate}
\end{theorem}

\begin{proof}
By simultaneous induction on the typing derivation
\(\Gamma \vdash t : A\).  We show the key cases:

\medskip\noindent
\textbf{Case (Fill).}
The typing rule (Fill) has premises:
\begin{itemize}
\item \(\forall \alpha.\; \Gamma \vdash_{\varphi_\alpha} t_\alpha :
  \restr{A}{\varphi_\alpha}\) --- by induction, \(\den{t_\alpha}\)
  is a local section of \(\den{A}\) on fiber \(\varphi_\alpha\).
\item \(\forall \alpha, \beta.\; \Gamma \vdash p_{\alpha\beta}\) ---
  by induction, the local sections are compatible.
\item Exhaustiveness: \(\vdash \bigvee_\alpha \varphi_\alpha\).
\end{itemize}
By Definition~\ref{def:descent-denotation}, the denotation
\(\den{\mathsf{fill}(\covers, \{t_\alpha\}, \{p_{\alpha\beta}\})}\)
is the unique global section of the sheaf \(\den{A}\).
Naturality holds because restriction commutes with the unique
global section construction.

\medskip\noindent
\textbf{Case (Res).}
The denotation
\(\den{\restr{t}{\varphi}}(\psi) = \den{t}(\psi \wedge \varphi)\)
is a section of \(\restr{\den{A}}{\varphi}\) by functoriality of
the restriction.

\medskip\noindent
\textbf{Case (HComp).}
Homogeneous composition is exactly descent on the interval face system.
The face predicates are \([i = 0]\) and \([i = 1]\) and the \(\psi_k\),
forming a cover of the interval.  The Kan condition (which holds
for our cubical sets) gives the filler, identical to the sheaf condition
on the covering.

\medskip\noindent
\textbf{Case (RefinI).}
\(\den{\langle a, h \rangle_\varphi}(\psi) =
(\den{a}(\psi \wedge \varphi), \den{h}(\psi \wedge \varphi))\)
is a section of \(\den{\{x : A \mid \varphi(x)\}}(\psi)
= \den{A}(\psi \wedge \varphi)\).

All other cases follow by standard induction arguments for CIC
and cubical type theories.  \hfill\(\square\)
\end{proof}

\begin{corollary}[Logical consistency]
\(\mathrm{C}^4\) is logically consistent: there is no closed term
\(t : \mathsf{False}\).
\end{corollary}

\begin{proof}
The denotation \(\den{\mathsf{False}}(\mathsf{True}) = \emptyset\)
(the empty set).  A closed term \(t : \mathsf{False}\) would have
\(\den{t} \in \emptyset\), which is impossible.
\end{proof}


% ─────────────────────────────────────────────────────────
\section{Adequacy: Operational Correspondence}
\label{sec:adequacy}

\begin{theorem}[Adequacy of the Computation Rules]
\label{thm:adequacy}
If \(\Gamma \vdash t : A\) and \(t \rightsquigarrow t'\)
via any computation rule of \(\mathrm{C}^4\), then
\(\den{t} = \den{t'}\) in \(\den{A}(\mathsf{True})\).
\end{theorem}

\begin{proof}
We check each computation rule:

\medskip\noindent
\textbf{(Fill-\(\beta\)):}
\(\restr{\mathsf{fill}(\covers, \{t_\alpha\}, \{p_{\alpha\beta}\})}{\varphi_k}
\rightsquigarrow t_k\).

By Definition~\ref{def:descent-denotation},
\(\restr{\den{\mathsf{fill}(\cdots)}}{\varphi_k}
= \den{t_k}\) (the sheaf section restricts to the local section).
So \(\den{\restr{\mathsf{fill}(\cdots)}{\varphi_k}} = \den{t_k}\).

\medskip\noindent
\textbf{(Res-\(\lambda\)):}
\(\restr{(\lambda x.\, t)}{\varphi}
\rightsquigarrow \lambda x.\, \restr{t}{\varphi}\).

By the denotation of \(\lambda\) and restriction:
\[
  \den{\restr{(\lambda x.\, t)}{\varphi}}(\psi) =
  \den{\lambda x.\, t}(\psi \wedge \varphi) =
  (a \mapsto \den{t}(\psi \wedge \varphi, x \mapsto a)) =
  \den{\lambda x.\, \restr{t}{\varphi}}(\psi).
\]

\medskip\noindent
\textbf{(\(\beta\)):}
\((\lambda(x : A).\, t)\; u \rightsquigarrow t[x := u]\).

Standard: the denotation of application is function application
in the model, which equals the denotation of the substitution.

\medskip\noindent
All other rules are verified similarly.  \hfill\(\square\)
\end{proof}

\begin{corollary}[Subject Reduction]
\label{cor:subject-reduction}
If \(\Gamma \vdash t : A\) and \(t \rightsquigarrow t'\), then
\(\Gamma \vdash t' : A\).
\end{corollary}

\begin{proof}
The reduction \(t \rightsquigarrow t'\) implies \(\den{t} = \den{t'}\)
(Theorem~\ref{thm:adequacy}).  By the soundness theorem
(Theorem~\ref{thm:soundness}), \(\den{t} \in \den{A}(\mathsf{True})\),
so \(\den{t'} \in \den{A}(\mathsf{True})\).  By adequacy of the
typing rules (the denotation is injective on closed normal forms,
proved by normalization), \(t'\) is well-typed at \(A\).
\end{proof}


% ─────────────────────────────────────────────────────────
\section{Trust Semantics}
\label{sec:trust-semantics}

The trust provenance tracked in judgment form~(6) has a precise
semantic interpretation.

\begin{definition}[Trust Sheaf]
The \textbf{trust sheaf} \(\mathcal{T}\) on \(\mathcal{C}_\alpha\) is the
constant sheaf valued in power sets:
\[
  \mathcal{T}(\varphi) \;\coloneqq\;
  \mathcal{P}(\{\mathsf{KERNEL}, \mathsf{Z3}, \mathsf{LEAN},
  \mathsf{LIBRARY}, \mathsf{ASSUMED}\}).
\]
Restriction maps are union: \(\restr{T}{f} = T\) (trust only grows).
\end{definition}

\begin{theorem}[Trust Monotonicity]
\label{thm:trust-mono}
Trust provenance is monotone under all \(\mathrm{C}^4\) rules:
\begin{enumerate}
\item \(\mathrm{trust}(\mathsf{fill}(\covers, \{t_\alpha\}, \{p\})) \supseteq
  \bigcup_\alpha \mathrm{trust}(t_\alpha)\).
\item \(\mathrm{trust}(\restr{t}{\varphi}) = \mathrm{trust}(t)\).
\item \(\mathrm{trust}(f\; u) \supseteq \mathrm{trust}(f) \cup \mathrm{trust}(u)\).
\item \(\mathrm{trust}(\transp^i\, A\, h\, t) \supseteq \mathrm{trust}(t)\).
\item \(\mathrm{trust}(\mathsf{import}(\mathit{name})) = \{\mathsf{LEAN}\}\).
\item \(\mathrm{trust}(\mathsf{assume}(P)) = \{\mathsf{ASSUMED}\}\).
\end{enumerate}
\end{theorem}

\begin{proof}
By the denotation of trust: the trust set of a term is the union of
the trust sets of all its sub-derivations.  Each rule only adds
sub-derivations, hence trust can only grow.  Formally, trust is
a section of the trust sheaf, and the trust sheaf's restriction maps
are identity (trust is a global constant on each fiber).
\end{proof}

\begin{corollary}[Trust as a Sheaf over the Cover]
For a descent proof \(\mathsf{fill}(\covers, \{t_\alpha\}, \{p\})\),
the trust of the global proof is the union of the trust of each
fiber proof:
\[
  \mathrm{trust}\bigl(\mathsf{fill}(\covers, \{t_\alpha\}, \{p\})\bigr)
  = \bigcup_\alpha \mathrm{trust}(t_\alpha)
  \cup \bigcup_{\alpha,\beta} \mathrm{trust}(p_{\alpha\beta}).
\]
This is the trust sheaf's global sections formula (Synergy~33).
\end{corollary}


% ─────────────────────────────────────────────────────────
\section{Comparison with the Lean Kernel Model}
\label{sec:lean-comparison}

The formal semantics of $\mathrm{C}^4$ can be compared with the
set-theoretic model of CIC used in the Lean~4 kernel:

\begin{center}
\begin{tabular}{lll}
\toprule
\textbf{Feature} & \textbf{Lean 4 model} & \textbf{$\mathrm{C}^4$ model} \\
\midrule
Base category & Trivial (1 object) & Refinement site \(\mathcal{C}_\alpha\) \\
Types & Sets & Sheaves over \(\mathcal{C}_\alpha\) \\
Terms & Set elements & Global sections \\
Equality & Propositional & Cubical paths \\
Decomposition & None (manual) & Descent rule \\
Decidability & Propositional (K axiom) & Z3 + Hedberg \\
Trust & None (all kernel) & Sheaf of provenance \\
\bottomrule
\end{tabular}
\end{center}

\(\mathrm{C}^4\) strictly extends Lean~4's model by adding:
\begin{enumerate}
\item The refinement site structure (giving decomposition along covers).
\item The cubical interval (giving proof-relevant equality and transport).
\item The trust sheaf (giving provenance tracking).
\end{enumerate}

The inclusion functor
\(\iota : \mathrm{Lean}_4 \to \mathrm{C}^4\) maps the trivial
site model to the \(\mathrm{C}^4\) model and is fully faithful on the
CIC fragment.  This is the semantic content of the conservativity theorem.


% ─────────────────────────────────────────────────────────
\section{Relation to the Mechanical Soundness Proof}
\label{sec:lean-proof-relation}

Chapter~\ref{ch:c4-lean} presents a mechanised soundness proof in
Lean~4, covering the following theorems from this chapter:

\begin{enumerate}
\item \textbf{Descent soundness}
  (Proposition from Definition~\ref{def:descent-denotation}):
  proved as \texttt{C4.Descent.soundness} in Lean.
\item \textbf{Subject reduction} (Corollary~\ref{cor:subject-reduction}):
  proved for the novel rules (descent-\(\beta\), restriction
  distribution) as \texttt{C4.SubjectReduction.sr\_fill\_beta}
  and related lemmas.
\item \textbf{Trust monotonicity}
  (Theorem~\ref{thm:trust-mono}):
  proved as \texttt{C4.Trust.monotone} in Lean.
\item \textbf{Hedberg collapse}:
  proved as \texttt{C4.Hedberg.propSheafTrivial} in Lean.
\item \textbf{Logical consistency}:
  proved as \texttt{C4.Soundness.consistency} in Lean.
\end{enumerate}

The Lean proofs mechanize the paper-level arguments given in this
chapter.  Where the paper relies on established results (CIC
normalization, cubical Kan condition), the Lean development uses
corresponding \texttt{axiom} declarations with references to the
published proofs.
